% Natural case studies: the template of the restructured manuscript's natural section (restructure plan, part 7), fixed
% before any analysis of natural-pr-v1 exists, so that its wording presupposes no outcome of the run.
% * Every value is a macro: \nrv{key} of generated/natural_run/numbers.tex (paper/make_tables.py --natural-run; ledger
%   paper/evidence_ledger_natural_run.md) or \nat{key} of generated/natural/numbers.tex (--natural; ledger
%   paper/evidence_ledger_natural.md). A key that is not generated yet prints UNFILLED.
% * Every statement cites the decision-log entry or the protocol section it comes from.
% * Premises: the fixed sentences state construction facts that the run's analysis also computes (the run's families,
%   invocations, revision and pilot flag; B1 against the regression norm; the labelers; the reference checks' fix
%   rounds; the failed audits). tests/test_natural_template.py (PREMISES) checks them on the synthetic world and, once it
%   exists, on the run's own numbers.tex. A premise the run does not bear out means that the sentence is rewritten from
%   the records, never that a cell is changed. The fix rounds are read from each reference-check status record's integer
%   fix_rounds, a key D64 does not fix: the build session confirms that the six records hold it before the run is
%   analysed, since the published analysis is a byte copy (scripts/publish_natural_run.py).
% * Families are named, never counted; no thread text, repository name, login or agent product is written.
% * main.tex inputs this file as the manuscript's natural section (the restructure). Its only pointers outside itself,
%   its generated tables and its part in App. G are the two that D77 left open, re-added once the appendix moves gave
%   them their content (App. A's generated table of amendments, its v0.12 row, and App. G's comparison of the natural arm
%   with the controlled study), and the controlled study's coverage table, to which the reading with the gate points.
% * Its part in App. G is paper/sections/app_natural_cases.tex (D80): what was moved there unchanged, so that the
%   main-text section stays short. Every rule of this header holds for it, and tests/test_natural_template.py reads both.
% * What changed after the run's analysis (D79): (1) the closing paragraph's integration slot, replaced by the one
%   sentence its rule asks for, the only text of the template written from the run's cells (the abstract's, the
%   introduction's and the discussion's sentences on the case studies, outside the template, were written after the
%   analysis too; D80); (2) the two pointers above; (3) review corrections of statements about the records, none of
%   which states an outcome of the run: what was fixed before the analysis and what the build session had seen (Status),
%   the caveat on D70's construction facts and the wording of the task-list exposure (Caveats), NP-05's correction as a
%   builder default and NP-20's near-miss disclosure as D77 records it (their paragraphs), and the planned cost per
%   accepted task, not reported (closing paragraph).
% * What changed in the next round after the analysis (D80): (4) the slot, rewritten with the gate's marker that Table
%   tab:nr-source now carries, by D77's rule (one sentence per pattern of the C1 and C3 cells, each family named once);
%   (5) the new paragraph that reads the families with the reference gate, and the pointer to the coverage table; (6) the
%   Status paragraph's last sentences; (7) moved unchanged to App. G (app_natural_cases.tex): the caveats on the build's
%   process, Table tab:nr-caveats, the sentence that names the parts of a family paragraph, the family paragraphs and the
%   cost per family, with new sentences that point to them; (8) where the records contradict the text: the closing
%   paragraph's "no statement here counts families", which its own screening count of S10 contradicts, now reads "by
%   their outcomes in the run"; and an added sentence after the caveat on the generators' training data, which the
%   paper cannot check.
% * What changed in a later round (D82): (9) moved unchanged to App. G (app_natural_cases.tex), so that the main text
%   carries no run bookkeeping: the Status paragraph's sentence on the run's start and end times, invocations and code
%   revision, with D76's account of the stopped first invocation, now the paragraph "The run's dates and invocations."
%   there; a new sentence in its place points to it, and App. G's first paragraph names it.
\providecommand{\nrcaveat}[1]{\raggedright\hspace{0pt}#1}

\section{Natural case studies}
\label{sec:natural-cases}

\paragraph{Status.}
The run \texttt{\NATURALRUN} graded the \nrv{families} accepted families of the first sample (\nrv{run.families}) under
harness \nrv{run.harness}: each family's frozen verifiers (C1 to C4, B1, B2, B3 and the reference checks) and its
independent labeler on every held-out candidate, each pair \nrv{run.repeats} times (protocol~\S20.2; D76). Its start
and end times, invocations and code revision, as its manifest records them, are in Appendix~\ref{app:natural-cases}
(D82). It is an automated
pilot: provenance eligibility needs a human reviewer and stays pending, so each family's acceptance gate is pending, and
the manifest's pilot flag is \nrv{run.pilot} (D64, decision~9; D75; D76). Protocol v0.12 (protocol~\S20) gave each
sample a run of its own after phases P1 to P4 of the first sample's families had run and before any natural held-out
candidate existed (D71; Table~\ref{tab:amendments}, its v0.12 row). When this run started, the extension's run,
\texttt{natural-pr-v2}, had not been made: the extension's S10 step was stopped at a usage limit, and the task
directories it left are to be built again from the start (D66, D75), with the results of this run known
(protocol~\S20.3). This section's wording and the run's analysis fields were fixed before any analysis of the run, while
grading was still in progress; before then, the build session's progress checks had shown it counts of outcomes over all
graded pairs together, naming no verifier or candidate (D77). After the analysis, the opening of the closing paragraph
was written from the run's cells, and review corrected or added statements about the records in this section, none of
which states an outcome of the run (D79). The sentences on these case studies in the abstract, the introduction and the
discussion were also written after the analysis, from the run's cells. A later round rewrote that opening with the
reference gate, added the reading of the families with it, and moved the family paragraphs and other detail of this
section, unchanged, to Appendix~\ref{app:natural-cases} (D80).

\paragraph{What was built.}
Each family follows the recipe of D64 in phases P0 to P7, and the end of each of P1 to P7 is a committed record (D65,
D74). The generated verifiers are built and frozen in P2 and P3. The reference contract is drafted in P4, after them, so
no generator saw it, and before the reference checks and the independent labeler of P5. The held-out pool of P6,
alternatives and near misses with a workspace each, starts after every verifier of its family is frozen, and P7 runs the
mappers, the duplicate check and the audits (D64; D72; D75). A family accepted at S11 is accepted for building: it passes
S10 within the family cap (protocol~\S16.4). The natural arm differs from the controlled study in these respects
(Table~\ref{tab:natural-vs-controlled}, Appendix~\ref{app:natural}):
\begin{enumerate}
\item The correction is a reviewer's comment on an agent's pull request. The reviewer need not be the person who opened
  the issue: we treat the comment as the correction, a construct difference from an in-session user correction (D64,
  decision~1).
\item The initial request is the linked issue, the pre-correction state is the tree the reviewer commented on (the
  pre-review state), and the reference is the addressing commit, which is evidence about the correction rather than an
  oracle (protocol~\S15.6; D70).
\item A configuration-A agent drafts the reference contract in P4 from the rendered thread and the addressing diff, and
  also from the reference diff, which adds the part of the change made before the review (D64, deviation~(k)).
\item Labels: D19, whose reference checks and regression norm are graded repeatedly on each candidate, a label being
  ambiguous where the repeats disagree (D64, decision~7); the source-evidence label of the pre-review states, the label
  of the first estimand of protocol~\S15.6 (D64, decision~14); and the independent labeler, a second labeling
  (protocol~\S15.5).
\item The audit holds the reference to be valid (D24, D27), so the natural C2 and C4 are audited against an unreliable
  oracle (D70). B1 equals D19's regression norm in NP-03, NP-05 and NP-20, and in part in NP-19 (D64, decision~3). B3
  is a secondary arm (protocol~\S15.6).
\end{enumerate}

\paragraph{Caveats.}
These hold for every number of this section and of its tables, whose captions point here with the marker
$^{\ddagger}$.
\begin{itemize}
\item The reference is not an oracle, and the reference checks were validated on it: they are usable only where the
  reference passes them in every repeat of their validation and the development no-op fails them (D64; D70).
\item The construction facts that the family paragraphs take from D70, such as why a reference fails a verifier or what
  revisers relaxed, are observations of the build made by model agents and not reviewed by a person (D70).
\item Every label is a model proposal: D19's, the independent labeler's, and the source-evidence label, which rests on
  the reviewer's request, the addressing commit and the agreement of the model judges of S9 (protocol~\S15.6). No human
  label exists (D64, decision~8).
\item The independent labelers of NP-03 and NP-19 are unusable, so the second labeling leaves them out (D72, D75, D76).
\item A candidate that D19 labels ambiguous is left out of false acceptance and false rejection, and is queued for human
  review (D19; D76).
\item The run is an automated pilot, as the status paragraph says (D64, decision~9; D76).
\item Caveats on how the families were built hold here too and are listed in Appendix~\ref{app:natural-cases}:
  the run's first step, the agents' context, and the runs that worked outside the recipe.
\item The run grades \nrv{accepted} accepted tasks of its own, fewer than \nrv{threshold}, so its families are
  exploratory case studies, each reported on its own, in per-family tables only, with no pooled rate, interval or
  p-value (protocol~\S19.6, \S20.3). Nor do we give a total over the families, or pool anything with the controlled arm
  (protocol~\S2).
\item The results of this run will be known when the extension's families are built; this is disclosed with the results
  of \texttt{natural-pr-v2} (protocol~\S20.3).
\item The construct differs from an in-session correction, as the first item of the list above says (D64, decision~1).
\item The pull requests and their repositories predate the generators' training data, so public merged fixes may be
  memorised (protocol~\S15.6). The protocol states this order of dates; no model identifier is named, so the paper
  cannot check it.
\item NP-04 and NP-05 come from the same repository (D63; the committed index's repository ids).
\end{itemize}

\paragraph{Reading the tables and the family paragraphs.}
Table~\ref{tab:nr-caveats} gives each family's caveats as values of the run's analysis, Table~\ref{tab:nr-source}
which conditions accept each family's pre-review state under the source-evidence label, and Table~\ref{tab:nr-fa-fr}
false acceptance and false rejection per condition under D19. Below, false acceptance is written accepted/invalid and
false rejection rejected/valid, under D19's labels, each distinct graded patch once (protocol~\S15.7); the pre-review
state counts under its D19 label, since Table~\ref{tab:nr-fa-fr} leaves out only the reference and the no-op. The
secondary oracle-gated view keeps a verifier's cell where the verifier passes the family's reference in the run and
drops it otherwise, a $^{\dagger}$ in Table~\ref{tab:nr-fa-fr} (protocol~\S15.7). Table~\ref{tab:nr-source} marks
the same cells, Table~\ref{tab:nr-sanity} gives each condition's outcome on the reference and on the no-op, the
unchanged base, and Table~\ref{tab:nr-coverage} how much of the reference contract each recovered contract covers, how
many of its requirements contradict the reference contract, and how many it marks superseded.
The family paragraphs, which set each family's cells beside its construction facts, its second labeling, its pairs and
its own caveats, are in Appendix~\ref{app:natural-cases}, with Table~\ref{tab:nr-caveats}.

\IfFileExists{generated/natural_run/natural_run_source_evidence.tex}{\begin{table}[t]\centering\footnotesize
\caption{Acceptance of each family's pre-review state, the tree the reviewer commented on before the correction, by condition, under the source-evidence label: invalid with respect to the correction, on the strength of the reviewer's request, the addressing commit and the agreement of both S9 judges (protocol \S15.6), so an accepted state is a false acceptance under this labeling. \textbf{Source evidence, not a human label}, reported apart from D19; D19's labels and the second labeling are model proposals, and no human labels exist. Run \texttt{natural-pr-v1} in case-study mode, as exploratory case studies (protocol \S19.6, \S20.3): the run grades 6 accepted tasks of its own (its families accepted in the committed index), fewer than 10, and each run is reported on its own, so every number is per family. Nothing is pooled over the families: no pooled rate, no interval and no p-value. Below the conditions: the same state's D19 label (ambiguous where its reference checks pass it, \S15.6), its label in the second labeling with the independent labeler's outcome (a pass on a pre-review state is ambiguous, decision log D64; left out: the family's labeler is not usable), and the reference checks' outcome on it in the run. A rejection is any outcome other than a pass: one that is not a task failure (an infrastructure outcome, or an unresolved pair whose repeats disagree, protocol \S15.3) is named in brackets. Terms: --: the run did not grade the condition for the family; not graded: the run did not grade the state with that verifier; none: no label (no D19 proposal, or no label in the second labeling). $^{\ddagger}$Natural case study; caveats in Section~\ref{sec:natural-cases}.}\label{tab:nr-source}
\begin{tabular}{p{0.25\linewidth}p{0.093\linewidth}p{0.093\linewidth}p{0.093\linewidth}p{0.093\linewidth}p{0.093\linewidth}p{0.093\linewidth}}\toprule Condition & NP-03 & NP-04 & NP-05 & NP-18 & NP-19 & NP-20 \\\midrule
\raggedright\hspace{0pt}C1 initial, plain & \raggedright\hspace{0pt}rejected & \raggedright\hspace{0pt}rejected & \raggedright\hspace{0pt}rejected & \raggedright\hspace{0pt}rejected & \raggedright\hspace{0pt}rejected & \raggedright\hspace{0pt}rejected \tabularnewline
\raggedright\hspace{0pt}C2 initial, audited & \raggedright\hspace{0pt}rejected & \raggedright\hspace{0pt}rejected & \raggedright\hspace{0pt}rejected & \raggedright\hspace{0pt}rejected & \raggedright\hspace{0pt}rejected & \raggedright\hspace{0pt}rejected \tabularnewline
\raggedright\hspace{0pt}C3 full, plain & \raggedright\hspace{0pt}rejected & \raggedright\hspace{0pt}rejected & \raggedright\hspace{0pt}rejected & \raggedright\hspace{0pt}rejected & \raggedright\hspace{0pt}rejected & \raggedright\hspace{0pt}rejected \tabularnewline
\raggedright\hspace{0pt}C4 full, audited & \raggedright\hspace{0pt}rejected & \raggedright\hspace{0pt}rejected & \raggedright\hspace{0pt}rejected & \raggedright\hspace{0pt}rejected & \raggedright\hspace{0pt}rejected & \raggedright\hspace{0pt}rejected \tabularnewline
\raggedright\hspace{0pt}B2 full-session LLM & \raggedright\hspace{0pt}rejected & \raggedright\hspace{0pt}rejected & \raggedright\hspace{0pt}rejected & \raggedright\hspace{0pt}rejected & \raggedright\hspace{0pt}rejected & \raggedright\hspace{0pt}rejected \tabularnewline
\raggedright\hspace{0pt}B1 repository tests & \raggedright\hspace{0pt}accepted & \raggedright\hspace{0pt}rejected & \raggedright\hspace{0pt}accepted & \raggedright\hspace{0pt}rejected & \raggedright\hspace{0pt}accepted & \raggedright\hspace{0pt}accepted \tabularnewline\midrule
\raggedright\hspace{0pt}B3 diff-derived (secondary) & \raggedright\hspace{0pt}rejected & \raggedright\hspace{0pt}accepted & \raggedright\hspace{0pt}rejected & \raggedright\hspace{0pt}rejected & \raggedright\hspace{0pt}rejected & \raggedright\hspace{0pt}rejected \tabularnewline\midrule
\raggedright\hspace{0pt}Its D19 label & \raggedright\hspace{0pt}invalid & \raggedright\hspace{0pt}invalid & \raggedright\hspace{0pt}invalid & \raggedright\hspace{0pt}invalid & \raggedright\hspace{0pt}invalid & \raggedright\hspace{0pt}invalid \tabularnewline
\raggedright\hspace{0pt}Its second label (the labeler's outcome) & \raggedright\hspace{0pt}left out & \raggedright\hspace{0pt}invalid (task failure) & \raggedright\hspace{0pt}invalid (task failure) & \raggedright\hspace{0pt}invalid (task failure) & \raggedright\hspace{0pt}left out & \raggedright\hspace{0pt}invalid (task failure) \tabularnewline
\raggedright\hspace{0pt}The reference checks' outcome on it & \raggedright\hspace{0pt}task failure & \raggedright\hspace{0pt}task failure & \raggedright\hspace{0pt}task failure & \raggedright\hspace{0pt}task failure & \raggedright\hspace{0pt}task failure & \raggedright\hspace{0pt}task failure \tabularnewline
\bottomrule\end{tabular}\end{table}
}{\unfilled{table of the pre-review states (make\_tables.py -{}-natural-run)}}
\IfFileExists{generated/natural_run/natural_run_fa_fr.tex}{\begin{table}[t]\centering\scriptsize\setlength{\tabcolsep}{4pt}
\caption{False acceptance $\cdot$ false rejection per family and condition, as accepted/invalid $\cdot$ rejected/valid, under D19's model-proposed labels, each distinct graded patch once; the reference solution and the no-op are excluded. A rejection is any outcome other than a pass, so an infrastructure outcome (environment failure, verifier error or timeout) and an unresolved pair (repeats that disagree, protocol \S15.3) count as rejections. Below the conditions: the family's pairs graded in the run (every verifier and candidate, label support included), those unresolved, and those that ended in an infrastructure outcome. \textbf{All labels are model proposals; no human labels exist.} Run \texttt{natural-pr-v1} in case-study mode, as exploratory case studies (protocol \S19.6, \S20.3): the run grades 6 accepted tasks of its own (its families accepted in the committed index), fewer than 10, and each run is reported on its own, so every number is per family. Nothing is pooled over the families: no pooled rate, no interval and no p-value. B3 (diff-derived) is a secondary arm (protocol \S15.6). n/a: the family has no unit of that kind under D19. --: the run did not grade the condition for the family, or the analysis does not hold the value. none: no such pair. $^{\dagger}$: the verifier does not pass the family's reference solution in the run, so the secondary oracle-gated view (protocol \S15.7) drops the cell and counts its units as lost yield. $^{\ddagger}$Natural case study; caveats in Section~\ref{sec:natural-cases}.}\label{tab:nr-fa-fr}
\begin{tabular}{p{0.2\linewidth}p{0.112\linewidth}p{0.112\linewidth}p{0.112\linewidth}p{0.112\linewidth}p{0.112\linewidth}p{0.112\linewidth}}\toprule Condition & NP-03 & NP-04 & NP-05 & NP-18 & NP-19 & NP-20 \\\midrule
\raggedright\hspace{0pt}C1 initial-only, plain & \raggedright\hspace{0pt}0/3 $\cdot$ 2/2$^{\dagger}$ & \raggedright\hspace{0pt}3/4 $\cdot$ 0/1 & \raggedright\hspace{0pt}0/3 $\cdot$ 2/2$^{\dagger}$ & \raggedright\hspace{0pt}0/4 $\cdot$ 2/2$^{\dagger}$ & \raggedright\hspace{0pt}2/4 $\cdot$ 0/2$^{\dagger}$ & \raggedright\hspace{0pt}0/3 $\cdot$ 2/2$^{\dagger}$ \tabularnewline
\raggedright\hspace{0pt}C2 initial-only, audited & \raggedright\hspace{0pt}1/3 $\cdot$ 2/2 & \raggedright\hspace{0pt}3/4 $\cdot$ 0/1 & \raggedright\hspace{0pt}1/3 $\cdot$ 0/2 & \raggedright\hspace{0pt}0/4 $\cdot$ 0/2 & \raggedright\hspace{0pt}1/4 $\cdot$ 0/2 & \raggedright\hspace{0pt}1/3 $\cdot$ 0/2 \tabularnewline
\raggedright\hspace{0pt}C3 correction-aware, plain & \raggedright\hspace{0pt}0/3 $\cdot$ 2/2$^{\dagger}$ & \raggedright\hspace{0pt}2/4 $\cdot$ 0/1 & \raggedright\hspace{0pt}0/3 $\cdot$ 2/2$^{\dagger}$ & \raggedright\hspace{0pt}0/4 $\cdot$ 2/2$^{\dagger}$ & \raggedright\hspace{0pt}1/4 $\cdot$ 0/2$^{\dagger}$ & \raggedright\hspace{0pt}0/3 $\cdot$ 0/2 \tabularnewline
\raggedright\hspace{0pt}C4 correction-aware, audited & \raggedright\hspace{0pt}0/3 $\cdot$ 2/2$^{\dagger}$ & \raggedright\hspace{0pt}2/4 $\cdot$ 0/1 & \raggedright\hspace{0pt}0/3 $\cdot$ 0/2 & \raggedright\hspace{0pt}0/4 $\cdot$ 0/2 & \raggedright\hspace{0pt}1/4 $\cdot$ 0/2$^{\dagger}$ & \raggedright\hspace{0pt}0/3 $\cdot$ 0/2 \tabularnewline
\raggedright\hspace{0pt}B2 session-LLM baseline & \raggedright\hspace{0pt}0/3 $\cdot$ 0/2 & \raggedright\hspace{0pt}1/4 $\cdot$ 0/1 & \raggedright\hspace{0pt}0/3 $\cdot$ 0/2 & \raggedright\hspace{0pt}0/4 $\cdot$ 0/2 & \raggedright\hspace{0pt}1/4 $\cdot$ 0/2 & \raggedright\hspace{0pt}0/3 $\cdot$ 0/2 \tabularnewline
\raggedright\hspace{0pt}B1 repository tests & \raggedright\hspace{0pt}3/3 $\cdot$ 0/2 & \raggedright\hspace{0pt}3/4 $\cdot$ 0/1 & \raggedright\hspace{0pt}3/3 $\cdot$ 0/2 & \raggedright\hspace{0pt}1/4 $\cdot$ 0/2 & \raggedright\hspace{0pt}4/4 $\cdot$ 0/2 & \raggedright\hspace{0pt}3/3 $\cdot$ 0/2 \tabularnewline\midrule
\raggedright\hspace{0pt}B3 diff-derived (secondary) & \raggedright\hspace{0pt}1/3 $\cdot$ 0/2 & \raggedright\hspace{0pt}4/4 $\cdot$ 0/1 & \raggedright\hspace{0pt}0/3 $\cdot$ 0/2 & \raggedright\hspace{0pt}0/4 $\cdot$ 0/2 & \raggedright\hspace{0pt}2/4 $\cdot$ 0/2 & \raggedright\hspace{0pt}1/3 $\cdot$ 0/2 \tabularnewline\midrule
\raggedright\hspace{0pt}The run's pairs: graded, unresolved & \raggedright\hspace{0pt}72, 0 & \raggedright\hspace{0pt}72, 0 & \raggedright\hspace{0pt}72, 0 & \raggedright\hspace{0pt}72, 0 & \raggedright\hspace{0pt}72, 0 & \raggedright\hspace{0pt}72, 0 \tabularnewline
\raggedright\hspace{0pt}The run's pairs ending in an infrastructure outcome & \raggedright\hspace{0pt}none & \raggedright\hspace{0pt}none & \raggedright\hspace{0pt}none & \raggedright\hspace{0pt}none & \raggedright\hspace{0pt}none & \raggedright\hspace{0pt}none \tabularnewline
\bottomrule\end{tabular}\end{table}
}{\unfilled{table of false acceptance and false rejection (make\_tables.py -{}-natural-run)}}

\paragraph{The families read with the reference gate.}
% Written after the analysis, from the run's cells (D80): outside the template's
% pre-commitment. tests/test_natural_template.py (READING_CELLS, READING_ORDER) checks each family it names against the
% run's own numbers and the published analysis.
A rejection of the pre-review state tells that state from the reference only where the verifier passes the reference,
and from valid work only where it also passes the valid alternatives; an acceptance says more than the base's only where
the verifier rejects the no-op. In NP-03, C1, C3 and C4 reject every candidate, the reference and both valid
alternatives included, and so do C1 and C3 in NP-05 and NP-18, and C1 in NP-20: their rejections of the pre-review state
there carry no information about the correction. In NP-03, NP-05, NP-19 and NP-20, B1 passes every candidate, the no-op
included, so its acceptance of the pre-review state there is the base's. In NP-04, where C1 and C3 pass the reference,
C3 and C4 accept eval-nm-2 and eval-nm-3 (false acceptance \nrv{c3.np04.fa} and \nrv{c4.np04.fa}). The near miss
eval-nm-2 targets R5, a correction requirement, and the label rule labels it invalid for breaking R5 alone, which the
correction-aware contract leaves uncovered under both mappers, so the coverage gap accounts for this acceptance and for
no other. The near miss eval-nm-3 targets R1 (D72), and the label rule labels it invalid for breaking R1, R4 and R5, of
which the correction-aware contract covers R1 and R4 under both mappers; every condition accepts it. C1, whose contract
covers neither correction requirement, accepts every near miss (\nrv{c1.np04.fa}). In NP-19, C3, C4 and B2 accept
eval-nm-1 (\nrv{c3.np19.fa}, \nrv{c4.np19.fa} and \nrv{b2.np19.fa}), a near miss on the correction; it targets R4, which
neither recovered contract covers (D75), but the label rule labels it invalid for breaking R1 to R4, R3 among them, a
correction requirement that the correction-aware contract covers under both mappers. The correction-aware contract
covers both correction requirements in NP-03, NP-05 and NP-18, and all but R11 in NP-20, and in NP-05 and NP-18 the
initial-only contract covers both as well; in the controlled tasks, by contrast, every correction-aware contract covers
its whole reference contract (Table~\ref{tab:coverage}). In NP-20 the correction-aware contract marks a requirement
superseded, and in NP-05 and NP-20 requirements of both recovered contracts contradict the reference contract under both
mappers; in NP-03, NP-04, NP-18 and NP-19 none does (Table~\ref{tab:nr-coverage}). B2's false acceptance and false
rejection are no higher than C3's or C4's in NP-03, NP-04, NP-05, NP-18, NP-19 and NP-20: B2 passes both valid
alternatives where C3 rejects them (NP-03, NP-05 and NP-18) and where C4 does (NP-03), accepts fewer near misses than
both in NP-04 (\nrv{b2.np04.fa}), and has the same cells as both in NP-19 and NP-20. C3's false acceptance is below C1's
in NP-04 and NP-19, its false rejection is below C1's in NP-20, and their cells are equal in NP-03, NP-05 and NP-18. The
audit cuts both ways: C2 passes the valid alternatives that C1 rejects in NP-05, NP-18 and NP-20, and C4 those that C3
rejects in NP-05 and NP-18, and C2 rejects an invalid candidate that C1 accepts in NP-19 (eval-nm-1), but C2 accepts an
invalid candidate that C1 rejects in NP-03, NP-05 and NP-20. These comparisons favour verifiers that agree with the
reference: the reference checks were validated on the addressing commit, whose diff the drafter of the reference
contract read, and a reference may not meet the contract recovered from its own thread.

\paragraph{What the case studies can and cannot show.}
% INTEGRATION SLOT (restructure plan, part 7, item 6). Its rule, as the template committed it before any analysis of
% the run (D77): the marker "prints UNFILLED until, after make_tables.py --natural-run has written the run's cells, it is
% replaced by sentences that name families by their cells and never count them, in the form "In NP-xx and NP-yy C1
% accepted the pre-review state and C3 rejected it", one sentence for each pattern of Table 7's C1 and C3 cells that
% occurs, each family named once. Nothing is written from the cells before then."
% Written after make_tables.py --natural-run had written the run's cells (D79), and rewritten after the gate's marker
% was added to Table tab:nr-source (D80): its C1 and C3 cells (se.c1.<family>,
% se.c3.<family>, and the marker that og.c1.<family> and og.c3.<family> hold) form three patterns, so the slot is the three
% sentences below, each family named once. tests/test_natural_template.py (INTEGRATION) checks the cells they state
% against the run's own numbers.
In NP-03, NP-05, NP-18 and NP-19, C1 and C3 rejected the pre-review state and failed the family's
reference, so the oracle-gated view drops both cells. In NP-04, C1 and C3 rejected the pre-review state and passed the
reference. In NP-20, C1 and C3 rejected the pre-review state; C1 failed the reference, so the oracle-gated view drops
its cell, and C3 passed it.
Table~\ref{tab:nr-source} gives, per family, whether C1 and C3 accept the pre-review state, and the family paragraphs
set each family's cells beside its caveats; no statement here counts families by their outcomes in the run. S10 passed
\nat{s10.first.passed} of the \nat{s10.first.processed} families of the first sample that it processed (D59): a
screening count, not a prevalence.
These case studies cannot show how often anything they show happens, nor whether it holds beyond these families.
